\begin{textsample}{POS Dim 4 – 2009\_c5 – Score -1.00 – 2009\_c5/enade\_2009\_psico\_95367afe99.txt}  \label{ex:f4_pos_267}
Capítulo 5 Distribuição dos Conceitos 5.1 Panorama nacional da distribuição dos conceitos A Tabela 5.1 apresenta a quantidade e o percentual de cursos de PSICOLOGIA participantes do ENADE/2009 que se situaram em cada uma das faixas de conceitos, distribuídos de acordo com a Região em que se localizam os cursos das instituições de educação superior, mantenedoras dos cursos responsáveis pela oferta.
O Gráfico 5.1 apresenta a mesma informação sob a forma de diagrama de caixa para uma melhor visualização da dispersão dos conceitos obtidos.
No diagrama de caixa, a caixa propriamente dita fica limitada pelos quartis.
A barra dentro da caixa em negrito representa a mediana.
Linhas verticais unem os quartis aos valores extremos.
Se estes valores extremos se configurarem como valores atípicos, estes aparecem como pontos isolados e as linhas terminam no maior ( ou menor ) valor ainda não considerado atípico.
Eventuais valores atípicos aparecem como pequenos círculos com a Categoria Administrativa correspondente.
Valores considerados muito atípicos foram representados por estrelas.
Um diagrama de caixa para poucos dados ou com dados muito concentrados pode ter alguns pontos coincidentes.
Por exemplo, se mais de metade dos cursos de uma Região receberem conceito 5, o valor extremo superior, o quartil superior e a mediana coincidirão.
Se o conjunto só contém um dado, este aparece denotado como a barra da mediana.
Se o conjunto apresenta todos os conceitos iguais, a notação é a mesma, somente uma barra horizontal.
Observando -se os dados da Tabela 5.1, percebe -se que, em termos nacionais, predominaram os cursos que receberam conceito 3, correspondendo a 37,9 \% do total nacional.
Pouco mais de 1 \% dos cursos foram classificados com conceito 1.
Conceitos 2 e 4 foram alocados a, respectivamente, 8,8 \% e 15,9 \%.
O conceito máximo foi alcançado por 4,5 \%.
Dos cursos de PSICOLOGIA, 125 ficaram sem conceito, correspondendo a 31,6 \% do total.
Considerando -se separadamente as cinco Regiões brasileiras ( ver também Gráfico 5.1 ), observa -se que, da Região Norte, 29 cursos participaram, dos quais a maioria, 53,7 \%, não obteve conceito.
Nenhum curso da Região obteve o conceito máximo, nem tampouco o conceito mínimo.
O conceito modal foi o conceito 3, com 24,1 \% e, respectivamente, 17,9 \% e 6,9 \% obtiveram os conceitos 2 e 4.
Esta Região foi a que apresentou cursos com a pior distribuição de conceitos.
Na Região Nordeste, 67 cursos participaram, dos quais 20,9 \% obtiveram conceito modal 3.
Nesta Região, 4 cursos ( 6,0 \% ) obtiveram o conceito máximo.
Nenhum curso obteve o conceito mínimo e 27 cursos ( 40,3 \% ) ficaram sem conceito.
Na Região Sudeste, 176 cursos participaram, dos quais, 76 obtiveram o conceito 3, equivalentes a 43,2 \% do total regional.
O conceito máximo foi alcançado por 7 cursos ( 4,0 \% ) e 4 cursos ( 2,3 \% ) obtiveram conceito mínimo.
Ficaram sem conceito 47 cursos ( 26,7 \% ).
Como a concentração no conceito 3 foi grande e centralizada, cursos com conceitos abaixo e acima deste valor aparecem como valores atípicos.
Na Região Sul, 40 dos 93 cursos avaliados, equivalentes a 43,0 \% do total regional, classificaram -se com conceito 3.
O conceito máximo foi alcançado por 5 cursos ( 5,4 \% ), 1 curso ( 1,1 \% ) obteve o conceito mínimo e 26 ( 28,0 \% ) cursos ficaram sem conceito.
Esta Região, juntamente com a Região Nordeste, foi a que apresentou cursos com a melhor distribuição de conceitos.
Na Região Centro-Oeste, 31 cursos participaram, sendo que 41,9 \% ( 13 cursos ) obtiveram o conceito modal 3.
O conceito máximo foi alcançado por 2 cursos ( 6,5 \% ).
Nenhum curso obteve o conceito mínimo e 10 ( 32,3 \% ) ficaram sem conceito.
Como a concentração no conceito 3 foi grande e centralizada, semelhantemente ao que ocorreu na Região Sudeste, cursos com conceitos abaixo e acima deste valor aparecem como valores atípicos. 5.2 Conceitos por Categoria Administrativa e por Região A Tabela 5.2 apresenta a distribuição dos cursos participantes do ENADE/2009 de PSICOLOGIA, por Categoria Administrativa, de acordo com os conceitos por eles alcançados, segundo Região.
O Gráfico 5.2 apresenta a mesma informação sob a forma de diagrama de caixa para uma melhor visualização da dispersão dos conceitos obtidos.
Somente eventuais valores atípicos aparecem com o identificador da Região na qual o curso foi habilitado.
De acordo com as informações apresentadas, a grande maioria, 82,6 \%, é de cursos em instituições privadas, 10,9 \% dos cursos eram vinculados a instituições federais, 2,8 \% a instituições estaduais e os restantes 3,8 \% a instituições municipais.
As instituições estaduais foram as que apresentaram, em termos nacionais, maior proporção de cursos com conceito 5, equivalente a 27,3 \% do total dessa categoria, seguidas pelas federais com 25,6 \%.
Os conceitos modais para as instituições federais foram 3 e 4, com 18,6 \% dos cursos em cada conceito ( 8 cursos cada ).
Entre as instituições estaduais, o comportamento foi semelhante e os conceitos modais são 3 e 4 com 18,2 \% dos cursos em cada.
Apenas 2 cursos ( 4,7 \% ) de instituições federais receberam o conceito mínimo, mas 12 ( 27,9 \% ) ficaram sem conceito.
Para as estaduais, 2 cursos receberam o conceito mínimo ( 18,2 \% ) e também 2 ficaram sem conceito.
Entre as instituições municipais, 40,0 \% receberam conceito modal 3.
Entre as instituições municipais, nenhuma recebeu o conceito máximo, 1 curso ( 6,7 \% ) recebeu o conceito mínimo e 6 cursos ( equivalentes a 40,0 \% nesta categoria ) ficaram sem conceito.
Os cursos nas instituições privadas também tiveram como conceito modal o valor 3 com uma incidência de 41,0 \%.
O conceito máximo foi alcançado por 4 cursos ( 1,2 \% ), nenhum recebeu o conceito mínimo e 32,1 \% ficaram sem conceito.
Considerando -se a análise por Região observou -se que apenas nas Regiões Sudeste e Sul, houve representação de cursos em todos os tipos de instituições.
Na Região Norte, dos 4 cursos participantes em instituições federais, 1 ficou sem conceito e os outros 3 cursos ficaram com os conceitos 3 e 4, respectivamente 50,0 \% e 25,0 \%.
Nenhum curso de instituição estadual participou e o único curso de instituição municipal participante obteve conceito 3.
Nessa Região, dos 24 cursos de instituições privadas participantes, 14 ( 58,3 \% ) ficaram sem conceito, 5 ( 20,8 \% ) obtiveram conceito 2, 4 obtiveram conceito 3 ( 16,7 \% ) e 1 obteve o conceito 4 ( 4,2 \% ).
Na Região Nordeste, dos 67 cursos participantes, 27 ( 30,8 \% ) ficaram sem conceito, assim distribuídos ; 4 em instituições federais, 2 em estaduais e 21 em instituições privadas.
Nesta Região, na área de PSICOLOGIA, nenhum curso de instituição municipal participou.
Nenhum curso na Região alcançou o conceito mínimo.
O máximo foi alcançado por apenas 3 cursos de instituições federais e 1 em instituição estadual.
O conceito 4 foi aplicado a 4 cursos de instituições federais, 1 curso de instituição estadual e 9 cursos de instituições privadas.
O conceito 3 foi atribuído a 1 curso de instituição federal, 1 de estadual e a 12 cursos de instituições privadas, enquanto que o conceito 2 foi alcançado por 1 curso de instituições federais e 7 cursos de instituições privadas.
Na Região Sudeste, dos 176 cursos participantes 47 ficaram sem conceito, sendo 2 em instituições federais, 5 em instituições municipais e 40 em instituições privadas.
O conceito modal para cursos em instituições federais foi o conceito máximo ( 36,4 \% ), nas estaduais, o conceito modal foi o conceito mínimo ( 66,7 \% ) e nas privadas foi o conceito 3 ( 48,1 \% ).
Nesta Região, os 3 cursos com conceito em instituições municipais receberam cada um os conceitos 1, 3 e 4.
Na Região Sul, dos 93 cursos participantes, 26 ficaram sem conceito, sendo 1 em instituição federal, 1 em instituição municipal e 24 em instituições privadas.
O conceito modal para cursos em instituições privadas foi 3 ( 43,8 \% ).
O conceito mínimo foi obtido por apenas 1 curso em instituição federal e o conceito máximo foi obtido por 5 cursos, 2 em instituição federal, 2 em instituição estadual e o outro em uma instituição privada.
Os 4 cursos avaliados em instituição municipal obtiveram conceito 3.
Na Região Centro-Oeste, dos 31 cursos participantes, 10 ficaram sem conceito, sendo 4 em instituições federais 6 em instituições privadas.
Nenhum curso de instituição estadual participou e o único curso de instituição municipal participante obteve conceito 2.
O conceito modal para cursos em instituições federais e privadas foi o conceito 3 ( 40,0 \% e 45,0 \%, respectivamente ).
Nesta Região, nenhum curso ficou com conceito mínimo.
Somente cursos em instituições federais alcançaram o conceito máximo nesta Região, 20,0 \%. 5.3 Conceitos por Organização Acadêmica e por Região Na Tabela 5.3 encontra -se a distribuição dos conceitos atribuídos aos cursos participantes do ENADE/2009 de PSICOLOGIA, por Organização Acadêmica, segundo as Regiões.
O Gráfico 5.3 apresenta a mesma informação sob a forma de diagrama de caixa para uma melhor visualização da dispersão dos conceitos obtidos.
Somente eventuais valores atípicos apareceram com o identificador da Categoria Administrativa.
Os cursos desta área estão principalmente alocados em Universidades ( 47,7 \% ), depois em Faculdades ( 38,9 \% ) e em Centros Universitários ( 13,4 \% ).
Este ordenamento se repete nas Regiões Sudeste, Sul e Centro-Oeste.
Nas Regiões Norte e Nordeste, os cursos em Faculdades foram mais numerosos do que os de Universidades Não foram avaliados cursos em PSICOLOGIA nem em Centros Federais de Educação Tecnológica ( CEFET ), nem em Institutos Federais de Educação, Ciência e Tecnologia ( IFECT ).
Em todas os tipos de organização administrativa observados, o conceito modal foi 3 : 46,6 \% nas Universidades, 56,6 \% nos Centros Universitários e 22,8 \% nas Faculdades.
Poucos cursos receberam conceito mínimo, e só em Universidades, 2,6 \%.
Apenas cursos em universidades e Faculdades alcançaram o conceito máximo ( respectivamente 8,5 \% e 1,3 \% ).
Dos 396 cursos avaliados, 125 ficaram sem conceito sendo, 23 em universidades ( 12,2 \% ), 15 em centros universitários ( 28,3 \% ) e 87 em faculdades ( 56,5 \% ).
Considerando -se separadamente as cinco Regiões brasileiras, observa -se que, na Região Norte, a categoria modal para Universidades foi 3 com 60,0 \% dos cursos avaliados na Região.
Nos Centros Universitários e Faculdades, a moda ficou nos conceitos 2 e 3, com respectivamente, 40,0 \% e 50,0 \% dos cursos avaliados para os dois tipos de Organização Acadêmica.
Nenhum curso nesta Região foi classificado nos conceitos mínimo e máximo.
Nesta Região, dos 29 cursos avaliados, 15 ficaram sem conceito sendo, 1 em universidade e 14 em faculdades.
Na Região Nordeste, o conceito modal entre os cursos avaliados em Universidades foi o conceito 4 ( 41,2 \% ) e em Centros Universitários e Faculdades foi o conceito 3 ( 66,7 \% e 40,0 \%, respectivamente ).
Apenas 4 cursos em Universidades obtiveram conceito máximo ( 23,5 \% dos avaliados na Região neste tipo de organização ).
Nenhum curso nesta Região foi classificado com o conceito mínimo.
Nesta Região, dos 67 cursos avaliados, 27 ficaram sem conceito sendo, 6 em universidades, 1 em centro universitário e 20 em faculdades.
Na Região Sudeste, os cursos avaliados de todos os tipos de organizações acadêmicas na área de PSICOLOGIA obtiveram o conceito 3 como modal : 54,9 \% em Universidades, 85,7 \% em Centros Universitários e 50,0 \% em Faculdades.
Apenas 5 cursos em universidades obtiveram conceito máximo ( 6,1 \% dos avaliados ).
O conceito mínimo também só foi aplicado a 4 cursos em universidades ( 4,9 \% ).
Nesta Região, dos 176 cursos avaliados, 47 ficaram sem conceito sendo, 8 em universidades, 11 em centros universitários e 28 em faculdades.
Na Região Sul, os cursos avaliados de todos os tipos de organizações acadêmicas na área também obtiveram o conceito 3 como modal : 56,3 \% em Universidades, 83,3 \% em Centros Universitários e 61,5 \% em Faculdades.
Somente 5 cursos em Universidades alcançaram conceito máximo, 10,4 \% dos avaliados e apenas 1 curso, também em universidade, obteve o conceito mínimo, 1,9 \%.
Nesta Região, dos 93 cursos avaliados, 26 ficaram sem conceito sendo, 4 em universidades, 3 em centros universitários e 19 em faculdades.
Na Região Centro-Oeste, o conceito modal dos cursos avaliados em Universidades e Centros Universitários na área de PSICOLOGIA também foi 3, respectivamente 64,3 \% e 100,0 \% dos cursos.
Nas Faculdades, o valor modal foi o conceito 2 com 50,0 \% dos cursos avaliados.
Nenhum curso nesta Região foi classificado com conceito mínimo e o conceito máximo foi obtido por apenas 2 cursos ( 14,3 \% ) em Universidades.
Nesta Região, dos 31 cursos avaliados, 10 ficaram sem conceito sendo, 4 em universidades e 6 em faculdades.

% matched lemmas: 
\end{textsample}
